\documentclass{jsarticle}
\usepackage[dvipdfmx, final]{graphicx}
\usepackage{amsmath}
\usepackage{amssymb}
\usepackage{chemfig}
\begin{document}
\title{}
\author{}
\date{}
\maketitle

 このレポートから、対数関数の底が$\zeta(s)$で$\log_{\zeta(k)}1=0$
 となり、$\zeta(k)={{\pi^{k}} \over {k!}}$が一般相対性理論の多様体積分
 の解が$e^{x\log x} + e^{-x\log x} \ge e^{x\log x} - e^{-x\log x}$
 と、シュバルツシルト半径のブラックホールと同じ解の様相を見せる
 $\zeta(k)={{\pi^{0}} \over {0!}}$で、$\log_{\zeta(k)}1!=0!$
 で、特異点が$0!$と宇宙の始まりは、${{\pi^{k}}\over {k!}}$という
 宇宙の種があると、ゼータ関数が展開されて、宇宙が始まり、
 

\end{document}


